\documentclass[10pt]{beamer}

% --- Theme and Colors ---
\usetheme{metropolis}
\usepackage[utf8]{inputenc}
\usepackage{booktabs}
\usepackage{listings}
\usepackage{xcolor}
\usepackage{amsmath}

% --- Code Styling ---
% \lstset{
%     basicstyle=\tiny\ttfamily,
%     breaklines=true,
%     keywordstyle=\color{blue},
%     commentstyle=\color{gray},
%     stringstyle=\color{orange},
%     frame=single,
%     backgroundcolor=\color{gray!5}
% }

\lstset{
    language=Python,
    basicstyle=\tiny\ttfamily,
    keywordstyle=\color{blue},
    stringstyle=\color{red},
    commentstyle=\color{gray},
    breaklines=true,
    frame=single,
    showstringspaces=false
}

\title{Midterm Progress:\\
Benchmarking LLMs on Edge Hardware}
\subtitle{Quantization Profiling and Runtime Optimizations}
\author{Group 5: Lau Wing Yi -- Bhandari, Om Raj}
\date{March 2026}

\begin{document}

\begin{frame}
    \titlepage
\end{frame}

% --- SECTION 1: SYSTEM DESIGN ---
\section{Preliminary System Design}

\begin{frame}{Progress: Automated Benchmarking Harness}
    Our \texttt{benchmark.py} suite automates the LLM inference benchmark lifecycle:
    \begin{enumerate}
        \item \textbf{Configuration:} Define device name, model and quantizations, context size (2048 by default), and batching (512 by default).
        \item \textbf{Environment Audit:} \texttt{print\_device\_info()} logs OS, CPU, and total RAM.
        \item \textbf{Acquisition:} \texttt{download\_models()} verifies GGUF local availability.
        \item \textbf{Profiling Loop:} \texttt{run\_quant\_benchmark()} for each precision level.
        \item \textbf{Ablation Study:} \texttt{run\_caching\_ablation()} for RAM vs. Disk caching.
        \item \textbf{Reporting:} Print matrix \& results
    \end{enumerate}
\end{frame}

% --- SECTION 2: EVALUATION ---
\section{Experimental Evaluation}

\begin{frame}{Preliminary Results: Benchmark Matrix}
    \textbf{Hardware:} x86 Linux Laptop (CPU-only, AMD Ryzen 7)
    \begin{table}[]
    \centering
    \resizebox{\textwidth}{!}{
    \begin{tabular}{lcccccc}
    \toprule
    \textbf{Quant} & \textbf{Load (s)} & \textbf{Peak RAM (MB)} & \textbf{TTFT (ms)} & \textbf{Prefill (tps)} & \textbf{Decode (tps)} \\ \midrule
    Q8\_0 & 0.35 & 8609.23 & 1836.46 & 55.54 & 9.99 \\
    Q4\_K\_M & 0.79 & 8356.86 & 1547.97 & 65.89 & 16.17 \\ \bottomrule
    \end{tabular}}
    \caption{Llama 3.1-8B on x86\_Arch\_Linux\_Laptop}
    \end{table}

    \textbf{Key Analysis:}
    \begin{itemize}
        \item \textbf{Q4 Load Time Paradox:} Longer load time ($0.79s$ vs $0.35s$) due to \emph{k-quants' complexity}. Unlike linear Q8, Q4\_K\_M requires a ``repacking" step to map its complex super-block structures to CPU SIMD instructions.
        \item \textbf{Q4 Memory Savings:} Q4 achieved the expected lower footprint, saving $\sim$253MB at peak.
    \end{itemize}
\end{frame}

\begin{frame}{The Mechanics of CPU Inference}
    \textbf{We expected Q8 to be faster than Q4 on a CPU, because:}
    \begin{itemize}
        \item \textbf{No Native 4-bit Opcode:} x86 ISAs (AVX2/AVX-512) lack 4-bit ALUs. Weights are packed in memory, then \textbf{unpacked} into 8-bit SIMD lanes at compute time.
    \end{itemize}
    \textbf{However, results show Q4 is faster. Why?}
    \begin{itemize}
        \item \textbf{The Bottleneck:} Modern CPUs are often \textbf{memory-bound}, not compute-bound.
        \item Smaller Q4 weights fit better in CPU caches and reduce memory bandwidth pressure. This performance gain outweighs the CPU cycles spent on the ``unpacking" logic.
    \end{itemize}
\end{frame}

\begin{frame}{Ablation Study: Prompt Caching}
    \begin{columns}
    \column{0.5\textwidth}
    \textbf{RAM Cache}
    \begin{itemize}
        \item Cold TTFT: $1990.45$ ms
        \item Warm TTFT: $1522.43$ ms
        \item Improvement: $468.02$ ms
        \item Reduction: $23.51\%$
    \end{itemize}
    \column{0.5\textwidth}
    \textbf{Disk Cache}
    \begin{itemize}
        \item Cold TTFT: $2582.98$ ms
        \item Warm TTFT: $1585.11$ ms
        \item Improvement: $997.87$ ms
        \item Reduction: $38.63\%$
    \end{itemize}
    \end{columns}
    
    \begin{itemize}
        \item \textbf{The OS Page Cache:} Disk ``warm" speed is identical to RAM. The OS stores the disk-based cache file in \textit{file-backed memory (RAM)} for quick access.
        \item \textbf{Strategic Choice:} 
        \begin{itemize}
            \item \textbf{RAM Cache:} Fast but volatile; takes up active address space.
            \item \textbf{Disk Cache:} Superior for edge devices. Provides near-RAM speeds via Page Cache (only if available!) but remains persistent and allows the OS to evict data under memory pressure.
            \item \textbf{Why it matters?} Some devices may have intermittent power!
        \end{itemize}
    \end{itemize}
\end{frame}

% --- SECTION 3: SYSTEM CHALLENGES ---
\section{Addressing System Challenges}

\begin{frame}{Ensuring Benchmarking Integrity}
    \begin{itemize}
        \item \textbf{Run Isolation:} Using \texttt{gc.collect()} and \texttt{time.sleep(2)} to clear memory fragmentation between tests in the same script instance.
        \item \textbf{The Warmup Protocol:} Implementing a \textit{warmup prompt} to force model weights into the page cache, preventing "Lazy Initialization" from skewing the first model's results.
        \item \textbf{Phase Isolation:} Using \textit{stream mode} to capture the exact timestamp of the first generated chunk, separating \textit{prefill} ($\approx$ TTFT) from \textit{decode} (TPS).
        \item \textbf{KV Cache Isolation:} A/B testing with fresh LLM instances to ensure we measure the cache effect specifically, not just context growth.
    \end{itemize}
\end{frame}


\begin{frame}[fragile]{Implementation: Profiling Core}
    \begin{lstlisting}[caption=Quantization Profiling Logic (SNIPPET)]
def run_quant_benchmark(model_path, quant):
    # 1. Baseline Memory Tracking
    process = psutil.Process()
    baseline_mem = process.memory_info().rss / (1024 * 1024)

    # 2. Lazy Loading / Model Init
    llm = Llama(model_path=model_path, n_ctx=CONTEXT_SIZE, verbose=False)
    
    # 3. CRITICAL: Warmup to ensure weights are in Page Cache
    _ = llm("Warmup", max_tokens=5)

    # 4. Phase Separation via Streaming
    start_eval = time.perf_counter()
    streamer = llm(prompt, max_tokens=100, stream=True)

    for chunk in streamer:
        if ttft_s == 0:
            # Captures Time to First Token (Prefill Phase)
            ttft_s = time.perf_counter() - start_eval
            decode_start_time = time.perf_counter()
        tokens_generated += 1

    # 5. Metrics Calculation (TPS and RAM Increase)
    decode_tps = (tokens_generated - 1) / (time.perf_counter() - decode_start_time)
    mem_increase = (process.memory_info().rss / (1024 * 1024)) - baseline_mem
    \end{lstlisting}
\end{frame}

% --- SECTION 4: FUTURE PLANS ---
\section{Future Plans}

\begin{frame}{Roadmap to Final Evaluation}
    \begin{itemize}
        \item \textbf{Test on more hardware:} Deploy to \textbf{NVIDIA Jetson X} and \textbf{MacBook (Metal optimization)}
        \item \textbf{Scaling \& Stress Testing:} Test long prompt sequences (P1 $\to$ P2 $\to$ P3) to identify the "breaking point" where RAM Cache hits capacity and requires swap space or Disk Caching.
        \item \textbf{Statistical Rigor:} Do more iterations to compute the averages with Standard Deviation.
    \end{itemize}
\end{frame}

\end{document}